\begin{table}[htbp]
\centering
\caption{Bid Aggressiveness: Price Effects Conditional on Competition}
\label{tab:bid_aggressiveness}
\begin{adjustbox}{max width=\textwidth}
\begin{threeparttable}
\small
\begin{tabular}{lcccccc}
\toprule
 & \multicolumn{3}{c}{Sequential controls} & \multicolumn{3}{c}{Fixed number of firms} \\
\cmidrule(lr){2-4} \cmidrule(lr){5-7}
 & (1) & (2) & (3) & (4) & (5) & (6) \\
 & Baseline & +Log firms & +SME winner & N=2 & N=3 & N$\geq$5 \\
\midrule
$g65 \times Pre$ & -0.1087*** & -0.1186*** & -0.1146*** & -0.0926*** & -0.1076*** & -0.1073*** \\
 & (0.0120) & (0.0123) & (0.0123) & (0.0166) & (0.0176) & (0.0206) \\
\midrule
Observations & 649,714 & 648,616 & 648,616 & 105,957 & 102,831 & 260,530 \\
R-squared & 0.9142 & 0.9148 & 0.9149 & 0.9173 & 0.9176 & 0.9299 \\
Item Fixed Effects & YES & YES & YES & YES & YES & YES \\
Month Fixed Effects & YES & YES & YES & YES & YES & YES \\
\bottomrule
\end{tabular}
\begin{tablenotes}
\small
\item \textit{Notes:} 18-month window, completed items. Standard errors clustered at the item level in parentheses.
Columns (1)--(3) sequentially add competition mediators to the baseline price regression.
Columns (4)--(6) restrict the sample to items with exactly 2, exactly 3, or 5+ participating firms,
holding the extensive margin approximately constant. The persistence of the price effect across all
subsamples rules out a pure headcount explanation for the cost reduction under open tenders; it does
not, on its own, identify the intensive margin structurally. Firm count is itself an equilibrium
outcome, so ``same N in both regimes'' does not guarantee the same bidder identities or cost
distribution.
*** \textit{p}$<$0.01, ** \textit{p}$<$0.05, * \textit{p}$<$0.1.
\end{tablenotes}
\end{threeparttable}
\end{adjustbox}
\end{table}
